% A1_STANDING
\subsection{Standing model and notation}

Let \(\epsilon_1=\epsilon_2=-1\) and
\(\epsilon_3=\epsilon_4=1\).  For \(a\in(-1,1)\), set
\[
 w_i(a)=\frac{1+\epsilon_i a}{4},\qquad
 s_0(a)=0,\qquad s_i(a)=\sum_{j\le i}w_j(a).
\]
Thus \(w_i(a)>0\), \(\sum_iw_i(a)=1\), and the signed one-step
current has mean \(a\).  Put
\[
 I_i(a)=(s_{i-1}(a),s_i(a)),\qquad
 J_i(a)=(s_{i-1}(a),s_i(a)).
\]
The seam-free four-branch baker correspondence is
\[
 B_a(q,p)=\left(
 \frac{q-s_{i-1}(a)}{w_i(a)},
 s_{i-1}(a)+w_i(a)p
 \right),\qquad q\in I_i(a),\quad 0<p<1.
\]
It maps the disjoint vertical ports \(I_i(a)\times(0,1)\) bijectively onto
the disjoint horizontal ports \((0,1)\times J_i(a)\).  On
\(\Sigma=\{1,2,3,4\}^{\mathbb Z}\), define
\[
 A_n(\omega)=\sum_{k=0}^{n-1}\epsilon_{\omega_k}.
\]
For a fixed reference \(a_0\),
\[
 \Lambda_{a_0}(\theta)
 =\log\sum_{i=1}^4w_i(a_0)e^{\theta\epsilon_i}
 =\log\big(\cosh\theta+a_0\sinh\theta\big),
\]
and its tilted mean is
\[
 a(\theta)=\Lambda_{a_0}'(\theta)
 =\frac{a_0+\tanh\theta}{1+a_0\tanh\theta}.
\]
All uses of \(w_i,s_i,B_a,A_n,\Lambda\), and the canonical field below refer
to these definitions.

% A1_MECHANICAL
\subsection{Autonomous Hamiltonian suspension of the cut-port map}

The mechanical section is a cut symplectic surface, not a torus.  Put
\[
 \mathcal Q_a^-
 =\bigsqcup_{i=1}^4 I_i(a)\times(0,1),
 \qquad
 \mathcal Q_a^+
 =\bigsqcup_{i=1}^4(0,1)\times J_i(a),
\]
with symplectic form \(\omega_0=dq\wedge dp\).  The closures are manifolds
with corners; all seam and corner faces have zero symplectic area.  On the
\(i\)-th component define
\[
 R_{a,i}(q,p)=
 \left(Q,P\right)
 =\left(
 \frac{q-s_{i-1}(a)}{w_i(a)},
 s_{i-1}(a)+w_i(a)p
 \right).
\]

\begin{lemma}[Exact cut-port symplecticity]
\label{lem:r4-a1-exact-port}
The disjoint-union map
\(R_a=\bigsqcup_iR_{a,i}:\mathcal Q_a^-\to\mathcal Q_a^+\) is an exact
symplectic bijection.  More precisely,
\[
 R_a^*(dQ\wedge dP)=dq\wedge dp,
\qquad
 R_a^*(P\,dQ)-p\,dq=dG_{a,i}
\]
on port \(i\), where
\[
 G_{a,i}(q,p)=\frac{s_{i-1}(a)}{w_i(a)}q.
\]
Its inverse is smooth on every outgoing port and the only loss of global
smoothness is the explicitly retained port boundary.
\end{lemma}

\begin{proof}
On port \(i\), \(dQ=w_i^{-1}dq\) and \(dP=w_i\,dp\), hence
\(dQ\wedge dP=dq\wedge dp\).  Moreover
\[
 P\,dQ-p\,dq
 =\left(s_{i-1}+w_i p\right)w_i^{-1}dq-p\,dq
 =\frac{s_{i-1}}{w_i}dq=dG_{a,i}.
\]
The inverse formulas
\(q=s_{i-1}+w_iQ\), \(p=(P-s_{i-1})/w_i\) prove the last assertion.
\end{proof}

Let
\[
 \widetilde{\mathcal M}_a
 =\mathcal Q_a^-\times[0,1]\times\mathbb R_E
\]
with two-form
\[
 \Omega_a=\omega_0+d\tau\wedge dE.
\]
Form the mapping torus
\[
 \mathcal M_a=
 \widetilde{\mathcal M}_a/\!\sim,
 \qquad
 (x,1,E)\sim(R_ax,0,E).
\]
The incoming and outgoing port labels are part of the quotient atlas.  Since
\(R_a\) is symplectic, \(\Omega_a\) descends to a nondegenerate closed
two-form on the seam-free mapping torus and to the natural symplectic form on
its manifold-with-corners completion.

\begin{theorem}[Cut-port Hamiltonian suspension]
\label{thm:r4-a1-suspension}
On \((\mathcal M_a,\Omega_a)\), the globally defined autonomous Hamiltonian
\[
 H_a(x,\tau,E)=E
\]
has a complete seam-free flow whose first return from \(\tau=0\) to
\(\tau=0\) is exactly \(R_a\), with return time one.  The invariant flux
measure on the section is \(dq\,dp\), and the flux of port \(i\) is
\(w_i(a)\).

Let \(c(x)=\epsilon_i\) on incoming port \(i\), and let
\(\alpha=c(x)d\tau\) on the suspension.  For every regular orbit with
itinerary \(\omega\),
\[
 \int_0^n\alpha(X_{H_a})\,dt=A_n(\omega).
\]
Consequently the canonical tilt by \(\theta A_n\) of the \(a_0\) Liouville
path law is exactly the Liouville path law of the mechanically distinct
member with parameter \(a(\theta)=\Lambda_{a_0}'(\theta)\).
\end{theorem}

\begin{proof}
With the convention
\(\iota_{X_H}\Omega_a=-dH\), one has \(X_{H_a}=\partial_\tau\).  The vector
field is compatible with the quotient identification because the gluing map
is independent of \(\tau,E\).  Starting at \((x,0,E)\), the orbit reaches
\((x,1,E)\), which represents \((R_ax,0,E)\); this proves the exact return
map and return time.  Liouville volume is
\(\Omega_a^2/2=dq\,dp\,d\tau\,dE\), so contraction with
\(X_{H_a}\) induces \(dq\,dp\) on the section.  Integrating it over
\(I_i(a)\times(0,1)\) gives \(w_i(a)\).

On the \(k\)-th suspension flight the port label is \(\omega_k\), hence
\(\alpha(X_{H_a})=\epsilon_{\omega_k}\); integration over its unit time and
summation gives \(A_n\).  Finally
\[
 \frac{w_i(a_0)e^{\theta\epsilon_i}}
 {\sum_jw_j(a_0)e^{\theta\epsilon_j}}
 =w_i(a(\theta)),
\]
so the tilted cylinder weights coincide at every word length.  Equality on
the cylinder sigma-field proves equality of the complete path laws.
\end{proof}

\begin{theorem}[Symplectic collar regularization]
\label{thm:r4-a1-collar}
For every compact set \(K\Subset\mathcal Q_a^-\) and every integer \(r\),
there is a smooth autonomous Hamiltonian channel on a smooth symplectic
six-manifold whose return map equals \(R_a\) on \(K\) and whose return maps
converge in \(C^r\) on compact port subsets to the cut-port map as the collar
width tends to zero.  The regularization preserves the port fluxes and the
work integral in the limit.
\end{theorem}

\begin{proof}
The exact primitive in Lemma~\ref{lem:r4-a1-exact-port} gives a compactly
supported Hamiltonian isotopy between the identity and each branch after
embedding the two port rectangles in disjoint Darboux collars.  Suspend the
isotopy by adjoining a clock--energy pair.  Because the four collars are
disjoint, their Hamiltonians may be summed.  Choose cutoffs which are one on
\(K\) and supported away from the corner faces.  Standard Moser correction
matches the collar forms on overlaps and leaves the core return map unchanged.
As the collars approach the port boundaries, the omitted area tends to zero
and the return maps converge with all derivatives on compact subsets.  The
Liouville flux and the bounded work one-form converge by dominated
convergence.
\end{proof}

% A1_CALIBRATION
\subsection{Canonical-cocycle rigidity on the common path space}

Let \((\mathcal F_t)\) be the coordinate filtration of \(\Sigma\).  A
conditional certainty equivalent \(\mathfrak V_{s,t}\) is assumed monotone,
cash additive, conditionally law invariant, continuous from above and below,
and strongly time consistent.  Assume additionally that it is generated by
one strictly monotone \(C^2\) utility \(u\) and is non-affine on some bounded
nonconstant random variable.

\begin{lemma}[Conditional Kolmogorov--Nagumo classification]
\label{lem:r4-a1-kn}
Under the preceding assumptions there is one constant \(\beta\ne0\),
independent of \(s,t\), such that
\[
 \mathfrak V_{s,t}(F)
 =\beta^{-1}\log
 \mathbb E_{\mathbb Q}[e^{\beta F}\mid\mathcal F_s].
\]
The affine limiting case is conditional expectation.
\end{lemma}

\begin{proof}
Cash additivity of the certainty equivalent implies that for every constant
\(c\) there are scalars \(A(c)>0,B(c)\) with
\[
 u(x+c)=A(c)u(x)+B(c).
\]
Composition of translations gives
\(A(c+d)=A(c)A(d)\); continuity yields \(A(c)=e^{\beta c}\).  If
\(\beta=0\), differentiating the translation equation gives an affine
utility.  Otherwise a second differentiation shows that
\(-u''/u'= -\beta\) is constant and hence, after an affine change of scale,
\(u(x)=e^{\beta x}\).  Conditional law invariance gives the conditional
quasi-arithmetic representation.  Strong time consistency applies the same
translation equation on every adjacent pair of time intervals and forces the
same \(\beta\) for all \(s,t\).
\end{proof}

\begin{theorem}[Mechanical calibration of the exponential coefficient]
\label{thm:r4-a1-calibration}
Suppose the reference probability \(\mathbb Q\) in
Lemma~\ref{lem:r4-a1-kn} is required to have, on every \(\mathcal F_n\), the
mechanical canonical cocycle
\[
 \frac{d\mathbb Q}{d\nu_{a_0}}
 =\exp\{\theta A_n-n\Lambda_{a_0}(\theta)\},
\]
and the same physical unit is used for \(A_n\) and for work in
\(\mathfrak V\).  Then \(\beta=\theta\) and
\[
 \mathfrak V_{s,t}(F)
 =\theta^{-1}
 \log\mathbb E_{\nu_{a(\theta)}}
 [e^{\theta F}\mid\mathcal F_s].
\]
Thus the coefficient is the mechanical conjugate field
\(I'(a)=\theta\), without assuming the exponential representation in the
calibration axiom.
\end{theorem}

\begin{proof}
Lemma~\ref{lem:r4-a1-kn} first derives the exponential representation and its
single coefficient \(\beta\).  Apply the resulting likelihood cocycle to the
one-step generator \(\epsilon_{\omega_0}\).  Its ratio between the positive
and negative symbol classes is \(e^{2\beta}\).  The mechanical cocycle has
ratio \(e^{2\theta}\).  Both classes have positive reference probability, so
equality of the two probability laws gives \(\beta=\theta\).  The unit
condition excludes replacing \(A_n\) by a scalar multiple.  The formula and
\(I'(a)=\theta\) follow from strict Legendre duality of
\(\Lambda_{a_0}\).
\end{proof}
